%% LyX 2.3.2 created this file.  For more info, see http://www.lyx.org/.
%% Do not edit unless you really know what you are doing.
\documentclass[11pt,american]{article}
\usepackage[T1]{fontenc}
\usepackage[utf8]{inputenc}
\usepackage[a4paper]{geometry}
\geometry{verbose,tmargin=0.6in,bmargin=0.6in,lmargin=0.7in,rmargin=0.7in}
\setlength{\parskip}{\medskipamount}
\setlength{\parindent}{0pt}
\usepackage{amssymb}
\usepackage{setspace}
\setstretch{0.9}
\usepackage{enumitem}

\makeatletter

%%%%%%%%%%%%%%%%%%%%%%%%%%%%%% LyX specific LaTeX commands.
\newcommand{\noun}[1]{\textsc{#1}}
%% Because html converters don't know tabularnewline
\providecommand{\tabularnewline}{\\}

%%%%%%%%%%%%%%%%%%%%%%%%%%%%%% Textclass specific LaTeX commands.
\newenvironment{lyxlist}[1]
	{\begin{list}{}
		{\settowidth{\labelwidth}{#1}
		 \setlength{\leftmargin}{\labelwidth}
		 \addtolength{\leftmargin}{\labelsep}
		 \renewcommand{\makelabel}[1]{##1\hfil}}}
	{\end{list}}

%%%%%%%%%%%%%%%%%%%%%%%%%%%%%% User specified LaTeX commands.
\frenchspacing
\newcommand{\bulletitem}{\raisebox{.3ex}{\tiny$\; \bullet \;$}}
\renewcommand{\arraystretch}{1.25}
% The original source loaded fancyheadings, but did not use it.

\usepackage{mathpazo}
\exhyphenpenalty=5000\hyphenpenalty=5000
\widowpenalty=10000
\clubpenalty=10000 
\makeatother

\usepackage[main=american,swedish]{babel}
\begin{document}
\begin{flushright}
\vspace{-12pt}
\textbf{\noun{\large{}Gustav Karreskog Rehbinder}}\vspace{-12pt}
\par\end{flushright}

\begin{flushright}
\emph{Telephone: }\textit{\emph{+46 (0) 762 10 42 20}}\emph{}\\
\emph{Email: }\textit{\emph{gustav.karreskog@gmail.com}}\\
\emph{Website: }\textit{\emph{ gustavkarreskog.com}}
\par\end{flushright}


\textbf{Current Position}\\
\rule[0.5ex]{1\linewidth}{1pt}
\vspace{-12pt}
\begin{lyxlist}{000.00.0000}
\item [{\noun{2025\textendash}}] \textbf{\emph{\noun{Researcher}}}\emph{
}\noun{}\\
Institute for Futures Studies, Stockholm\\
\end{lyxlist}

\textbf{Previous Positions}\\
\rule[0.5ex]{1\linewidth}{1pt}
\vspace{-12pt}
\begin{lyxlist}{000.00.0000}
\item [{\noun{2021\textendash2025}}] \textbf{\emph{\noun{Postdoctoral Researcher}}}\emph{
}\noun{}\\
Dept. of Economics, Uppsala University\\
\end{lyxlist}
\textbf{Education}\\
\rule[0.5ex]{1\linewidth}{1pt}

\vspace{-12pt}

\begin{lyxlist}{000.00.0000}
\item [{\noun{2015\textendash2021}}] \textbf{\emph{\noun{Doctor of Philosophy in Economics}}}\emph{
}\noun{}\\
Stockholm School of Economics\\
\emph{Primary supervisor (from 2021-01-01):} Assoc. Prof. Mark Voorneveld\\
\emph{Primary supervisor (until 2020-12-31):} Prof. Jörgen Weibull\\
\emph{Secondary supervisor:} Assoc. Prof. Erik Mohlin\\
\emph{\ }\\
\emph{References}\\
\begin{tabular}{lll}
\textbf{Prof. Jörgen Weibull} & \hspace{8bp} & \textbf{Prof. Drew Fudenberg}\tabularnewline
Dept. of Economics, SSE &  & Dept. of Economics, MIT\tabularnewline 
\textit{jorgen.weibull@hhs.se} &  & \textit{drewf@mit.edu}\tabularnewline
\textbf{Assoc. Prof. Erik Mohlin} &  & \textbf{Assoc. Prof. Mark Voorneveld}\tabularnewline
Dept. of Economics, Lund University &  & Dept. of Economics, SSE\tabularnewline
\textit{erik.mohlin@nek.lu.se} &  & \textit{mark.voorneveld@hhs.se}\tabularnewline
\end{tabular}
\item [{\noun{2014\textendash2015}}] \textbf{\emph{\noun{M.Sc. in Economics}}}\noun{}\\
Stockholm School of Economics\\
\emph{Unfinished due to admittance to the PhD program}
\item [{\noun{2012\textendash2014}}] \textbf{\emph{\noun{M.Sc. in Mathematics}}}\noun{}\\
Stockholm University
\item [{\noun{2010\textendash2012}}] \textbf{\emph{\noun{B.Sc. in Mathematics}}}\noun{}\\
Stockholm University
\end{lyxlist}
\vspace{-8pt}

\textbf{Research Visits}\\
\rule[0.5ex]{1\linewidth}{1pt}

\vspace{-12pt}

\begin{lyxlist}{000.00.0000}
\item [{\noun{2018\textendash2019}}] \textbf{{Dept. of Economics, Massachusetts Institute of Technology}}\noun{}\\
Faculty Sponsor: Prof. Drew Fudenberg
\end{lyxlist}
\vspace{-8pt}

\textbf{Teaching and Research Fields}\\
\rule[0.5ex]{1\linewidth}{1pt}

\vspace{-8pt}

\begin{lyxlist}{00.00.0000}
\item [{\noun{fields}}] \textit{\emph{Microeconomic Theory, Behavioral Economics, Experimental Economics}}
\item [{\noun{topics}}] \textit{\emph{Machine Learning, Bounded Rationality, Learning in Games}}
\end{lyxlist}
\vspace{-8pt}

\textbf{Journal Publications}\\
\rule[0.5ex]{1\linewidth}{1pt}

\vspace{2pt}

``Between Scylla and Charybdis: The Trade-Offs of Coalition Formation under Radical-Right Pressure'' \\
\emph{with Anders Kärnä, Jaakko Meriläinen and John Norell; Economica, forthcoming}

``Probabilistic functionalism as a limiting condition for robustness'' \\
\emph{with Mattias Forsgren and Benjamin Mandl; Scientific Reports 15, 44333 (2025)}

``Predicting Cooperation with Learning Models'' \\
\emph{with Drew Fudenberg; American Economic Journal: Microeconomics 16 (1), 1--32 (2024)}

\vspace{4pt}

\newpage

\textbf{Working Papers}\\
\rule[0.5ex]{1\linewidth}{1pt}

\vspace{2pt}

``Rational Heuristics for One-Shot Games'' \\
\emph{with Frederick Callaway and Thomas L. Griffiths}

``Stochastic Stability of a Recency Weighted Sampling Dynamic'' \\
\emph{with Alexander Aurell}

\textbf{Ongoing projects}\\
\rule[0.5ex]{1\linewidth}{1pt}

\vspace{2pt}

``Cue Based Decision Making and Context Effects'' \\
\emph{with Benjamin Mandl}

``Estimation of Learning Models Using Approximate Bayesian Computation''
\vspace{2pt}

\textbf{Journal Publication in Mathematics}\\
\rule[0.5ex]{1\linewidth}{1pt}

\vspace{2pt}

``Schrödinger operators on graphs: symmetrization and Eulerian cycles'' \\
\emph{Proceedings of the American Mathematical Society}, 144, (2016) \\
\emph{with Isak Trygg Kupersmidt and Pavel Kurasov}

\vspace{2pt}
\textbf{Research Grants and Awards}\\
\rule[0.5ex]{1\linewidth}{1pt}

\vspace{-8pt}

\begin{lyxlist}{MMMMM}
\item [{2017}] Tom Hedelius Scholarship for research visit to MIT. 
\item [{2014}] Scholarship for excellent Master Thesis from Mittag-Leffler's fund.
\end{lyxlist}

\vspace{-8pt}

\textbf{Teaching}\\
\rule[0.5ex]{1\linewidth}{1pt}

\vspace{-8pt}

% \begin{lyxlist}{00.00.0000}
% \item \textbf{Stockholm School of Economics}
% \item [{2020}] TA: Global Challenges -  Undergraduate course on conflict and cooperation
% \item [{2017}] TA: Economics of Organization - Undergraduate course 
% \item [{2016,2017}] Math summer camp - Preparatory math class for incoming Ph.D. students
% \item [{2016,2017}] TA: Mathematics I - Introductory mathematics for Ph.D. students
% \item [{2016}] TA: Advanced Microeconomics - Advanced level course on microeconomic theory
% \item [{2012\textendash 2015}] \textbf{Amanuensis}, Department of Mathematics, Stockholm University \\ 
% Primarily teaching assistant in undergraduate mathematics. I also developed (designed and coded) a web-platform  for a large distance course in preparatory mathematics.
% \end{lyxlist}

\begin{enumerate}[noitemsep]
\item[{}] \textbf{Uppsala University}
\item [{2022\textendash2023}] Lecturer: Nationalekonomi B: Mikroteori med tillämpningar (Intermediate Microeconomics)
\item [{2021\textendash2023}] Lecturer: Analytical Methods
\item[{}] \textbf{Stockholm School of Economics}
\item [{2020}] TA: Global Challenges -  Undergraduate course 
\item [{2017}] TA: Economics of Organization - Undergraduate course 
\item [{2016,2017}] Math summer camp - Preparatory math class for incoming Ph.D. students
\item [{2016,2017}] TA: Mathematics I - Introductory mathematics for Ph.D. students
\item [{2016}] TA: Advanced Microeconomics - Advanced level course on microeconomic theory
\item[{}]
\item [{2012\textendash 2015}] \textbf{Amanuensis}, Dept. of Mathematics, Stockholm University \\ 
Primarily teaching assistant in undergraduate mathematics. I also developed (designed and coded) a web-platform  for a large distance course in preparatory mathematics.
\end{enumerate}
\vspace{-5pt}

\textbf{Pedagogical Training}\\
\rule[0.5ex]{1\linewidth}{1pt}
\vspace{-8pt}

\begin{enumerate}[noitemsep]
\item [{2023}] Högskolepedagogisk grundkurs (five weeks), Uppsala University
\end{enumerate}
\vspace{-5pt}

\textbf{Presentations Outside of Home Department}\\
\rule[0.5ex]{1\linewidth}{1pt}

\vspace{-8pt}



\begin{enumerate}[noitemsep]
\item[{2021}] Arne Ryde Workshop, Lund; Games, 6th World Congress of the Game Theory Society, Budapest; Applied Micro Seminar, Goethe University Frankfurt 
\item[{2020}] Nordic Exchange, NHH, Bergen; SUDSWEC, Uppsala University; ENTER/SWIPS, University College London
\item[{2019}] Phd Math Fest, Stockholm; SING 15, Turku;  Theory Lunch, MIT, Boston
\item[{2018}] Theory Lunch, MIT, Boston
\end{enumerate}

\vspace{-5pt}

\textbf{Other Skills}\\
\rule[0.5ex]{1\linewidth}{1pt}
% \vspace{-10pt}
\textbf{Languages} \\ 
Swedish (native), English (fluent), Spanish (very good) \\

\textbf{Programming}\\
Julia, Python, R, Web Development (HTML, CSS, javascript, SQL, basic linux server administration etc.), and workable knowledge in many more such as STATA, Matlab, and Mathematica.



\pagebreak

\vspace{-10pt}
\textbf{Abstracts of selected papers}\\
\rule[0.5ex]{1\linewidth}{1pt}

\textbf{\underline{Working Paper}}\textbf{\textit{}}\\
\textbf{\textit{}}\\
\textbf{\textit{``Rational Heuristics for One-Shot Games'' }} \\ 
\emph{with Frederick Callaway and Thomas L. Griffiths (Dept. of Psychology, Princeton University)}


Insights from behavioral economics suggest that perfect rationality is an insufficient model of human decision-making. However, the empirically observed deviations from perfect rationality or biases vary substantially among environments. There is, therefore, a need for theories that inform us when and how we should expect deviations from rational behavior. We suggest that such a theory can be found by assuming optimal use of limited cognitive resources. In this paper, we present a theory of human behavior in one-shot interactions based on the rational use of heuristics. We test our theory by defining a broad family of heuristics for one-shot games and associated cognitive cost functions. In a large, preregistered experiment, we find that behavior is well predicted by our theory, which yields better predictions than existing models. We find that the participants' actions depend on their environment and previous experiences, in the way the rational use of heuristics suggest.

\textbf{\underline{Journal Publications}}\\
\textbf{\textit{``Between Scylla and Charybdis: The Trade-Offs of Coalition Formation under Radical-Right Pressure''}}\\
\emph{with Anders Kärnä, Jaakko Meriläinen and John Norell (Economica, forthcoming)}

The electoral success of populist radical-right parties makes non-cooperation increasingly costly for established parties. Using Swedish municipal politics, we examine when the Sweden Democrats become part of governing coalitions. We find a threshold of approximately 18.5\% electoral support, beyond which their inclusion becomes substantially more likely. Coalition size, office rents, and ideological dispersion change as their support approaches and crosses this threshold.

\textbf{\textit{``Probabilistic functionalism as a limiting condition for robustness''}}\\
\emph{with Mattias Forsgren and Benjamin Mandl (Scientific Reports, 2025)}

Behavioural phenomena often involve stimulus characteristics that may serve as cues to outcomes. People can learn whether such cues guide them towards their goals and adapt accordingly. In an experiment, both the attraction effect and the default nudge vary with how reliably those characteristics predict the superior option. Cue-outcome relationships may therefore limit the robustness of behavioural phenomena.

\textbf{\textit{``Predicting Cooperation with Learning Models''}}\\
\emph{with Drew Fudenberg (Dept. of Economics, MIT)}

We use simulations of a simple learning model to predict how cooperation varies with treatment  in the experimental play of the indefinitely repeated prisoner's dilemma.  We suppose that learning and the game parameters only influence play in the initial round of each supergame, and that after these rounds play depends only on the outcome of the previous round.  Using data from 17 papers, we find that our model predicts out-of-sample cooperation at least as well as more complicated models with more parameters  and harder-to-interpret  machine learning algorithms.  Our results let us predict how cooperation rates change with longer experimental sessions, and help explain  past findings on the role of strategic uncertainty.
\vspace{5mm}

\textbf{\underline{Working Papers}}\\
\textbf{\textit{``Stochastic Stability of a Recency Weighted Sampling Dynamic''}}\\
\emph{with Alexander Aurell (Dept. of Operations Research and Financial Engineering, Princeton University)}

We introduce and study a model of long-run convention formation for rare interactions. Players in this model form beliefs by observing a recency-weighted sample of past interactions, to which they noisily best respond. We propose a continuous state Markov model, well-suited for our setting, and develop a methodology that is relevant for a larger class of similar learning models. We show that the model admits a unique asymptotic distribution which concentrates its mass on some minimal CURB block configuration. In contrast to existing literature of long-run convention formation, we focus on behavior inside minimal CURB blocks and provide conditions for convergence to (approximate) mixed equilibria conventions inside minimal CURB blocks. 


\end{document}
